\exnumber{00}
\extitle{UX III - Ex 1 - Analyse concurrentielle}
\exfiles{ex00/archive zip, nomme Name\_FirstName\_DeliverableName\_Date\_VersionNumber}
\exauthorize{Ceux nommés dans l'exercice}
\makeheaderfiles

\textbf{Contexte :}

La quête n°3 porte sur l'Empathie.

Dans ce scénario, le principal de votre lycée a appris que vous étiez designer. Il a choisi de faire appel à vous pour l'aider à faire face au harcèlement scolaire parmi les adolescents.

Le harcèlement scolaire décrit des comportements de harcèlement en milieu scolaire. Il se caractérise par l'usage répété de violence, y compris la moquerie et d'autres humiliations.

En fait, ils se sont rendu compte qu'ils ne connaissaient pas très bien cette génération (15 à 18 ans). Ils vous demandent donc de mener une recherche utilisateur pour rencontrer cette population et la comprendre.

Votre travail est d'explorer votre sujet et d'utiliser divers outils inclus dans cette phase « Empathize ».

\textbf{Instructions :}

Menez une analyse concurrentielle parmi 10 associations ou organisations luttant contre le harcèlement scolaire.

Vous pouvez choisir 10 organisations dans le monde, ou de votre région ou de votre pays d'origine. Naviguons !

Comparez les actions principales, canaux et valeurs de ces 10 organisations, selon les critères que vous jugez pertinents.

\textbf{Livrables :}

\begin{itemize}
\item Un rapport concurrentiel avec une introduction qui explique vos conclusions principales (PDF)
\item Un graphique pour montrer le paysage concurrentiel, avec 2 axes de votre choix (numérique, psychologique, local, proximité des étudiants, institutionnel, etc.) (PDF)
\item Un tableau pour comparer les organisations sur 5 à 10 critères (PDF)
\item Déposez les documents sur Github, dans un zip nommé « ProjectTitle\_Name\_FirstName », et chaque livrable ainsi : « Name\_FirstName\_DeliverableName\_Date\_VersionNumber ». Exemple : « Doe\_John\_MidFiPrototype\_05242024\_V1 ».
\item Ajoutez un titre dans tous les documents écrits.
\item Soignez l'aspect global des livrables. Ils doivent être clairs, simples et lisibles. Vous pouvez regarder canva.com, sans surcharger le dessin de détails.
\end{itemize}

Ressources :

\begin{itemize}
\item Analyse concurrentielle : comment l'aborder efficacement \url{https://uxdesign.cc/competitive-analysis-how-to-approach-it-effectively-def97d130d2c}
\item Templates d'analyse concurrentielle \url{https://venngage.com/blog/competitor-analysis-template/}
\end{itemize}

\textbf{Citation :}

\begin{itemize}
\item « Concevez toujours une chose dans son contexte plus large : une chaise dans une pièce, une pièce dans une maison, une maison dans un environnement, un environnement dans un plan urbain » Eero Sarrinen | Knoll
\item Créer un grand design ne consiste pas seulement à comprendre ce que l'utilisateur veut. Il s'agit aussi de comprendre et de répondre aux besoins de l'entreprise.
\end{itemize}

